\addcontentsline{toc}{section}{Заключение}

В ходе проведённого исследования был осуществлён сравнительный анализ современных технологий разработки интеллектуальных систем: языков описания онтологий и функциональных языков программирования, а также инструментов для непрерывного извлечения знаний из интернета и англо-китайского билингвального графа знаний.

В первом разделе работы были подробно рассмотрены языки CycL (язык описания онтологий) и LISP (функциональный язык программирования). Проанализированы их синтаксические и семантические особенности, выявлены общие черты и различия. Результаты анализа были формализованы с помощью SCn-кода.

Во втором разделе исследованы современные инструменты для непрерывного извлечения знаний из интернета NELL и англо-китайский билингвальный граф знаний XLore. Проведен анализ их функциональных возможностей, а также особенностей архитектуры. Информация о каждом инструменте также была формализована на SCn-коде.

Результаты работы могут быть использованы при проектировании гибридных интеллектуальных систем, сочетающих возможности логического вывода и онтологического моделирования, а также при выборе и обосновании технологических решений для конкретных прикладных задач.